\section{Introduction}
\label{sec:intro}

Transformer inference increasingly uses formats narrower than the binary64
arithmetic for which much of classical numerical analysis was developed. The
dominant response is quantization: choose a target format, map weights and
activations onto it, and recover accuracy through calibration, outlier handling,
or per-channel scaling
\citep{3600270.3602468,xiao2023smoothquant,lin2024awq,gong2024what}. In that
program, the rounding operator is usually fixed by the underlying hardware format.
In post-training quantization, offline weight-rounding methods such as
AdaRound~\citep{nagel2020adaround} optimize the rounding direction of static
weights via a local second-order task loss surrogate. However, the rounding
rule used for intermediate calculations during inference has received less
attention. In contrast, dynamic rounding is already becoming explicit in
low-precision training: NVIDIA's microscaled NVFP4 recipe,
for example, selectively uses stochastic rounding for gradient casts to reduce
quantization bias~\citep{nvidia2025nvfp4}. Whether this same
selectivity is useful during inference remains an open question.

This paper isolates that question by holding the numerical format fixed and
varying only the rounding rule. Round-to-nearest~\citep{muller2018handbook} (RN)
returns the closer representable neighbor, whereas stochastic
rounding~\citep{forsythe1959rounding,croci2022stochastic} (SR), proposed by
Forsythe for controlling accumulated roundoff and pioneered in deep learning by
\citet{gupta2015deep} to prevent gradient stagnation in low-precision training,
selects either neighbor with probability proportional to proximity, making each
rounding error zero in expectation.
RN errors, by contrast, are deterministic functions of their operands and may
align; sufficiently small addends can be absorbed by a large accumulator or
repeated operations can drift away from the exact result. SR therefore trades
systematic local drift for higher variance.

Whether stochastic rounding improves inference accuracy depends first on the
operation. A long reduction is more exposed to accumulated RN drift than a short
one. In a GPT-2-style block~\citep{radford2019language}, the multilayer
perceptron (MLP) down-projection reduces over $d_{\rm ff}=3072$ terms, while an
attention-head score reduces over $d_{\rm head}=64$. At the same accumulator
precision, these operations can therefore have very different exposure to
accumulated roundoff.

The outcome also depends on the operation's location, both because local input
distributions affect how errors initially accumulate, and because downstream
layers transform the resulting perturbations. A perturbation in an attention score
passes through the softmax before changing the weighted value sum, whereas a
perturbation from the MLP down-projection enters the residual stream and is
transformed by later blocks. The same change of rounding rule can therefore
reduce perplexity at one site and increase it at another. This leads to the
research question of this paper: \emph{when and where does stochastic rounding
      improve transformer inference accuracy over round-to-nearest at low fixed
      precision?}

We study this question in the GPT-2 architecture, using DistilGPT-2 as a
six-block instance, and compare the two rounding rules site by site. Hardware
SR is tied to particular formats, operations, and accelerators, so experiments
at freely selected precisions require software emulation at every instrumented
operation, and the arithmetic intensity of transformer inference requires that
emulator to be fast and to vary precision and rounding rule at run time. The
contributions of this paper are as follows:

\begin{itemize}[leftmargin=1.5em]
      \item We present the \textbf{variable-precision stochastic rounding
                  (VPSR)} algorithm (\cref{sec:vpsr}), prove its correctness, and
            implement it in \textbf{PRISM}~\citep{prism}, where it runs up to
            $70\times$ faster than Verificarlo's Monte Carlo Arithmetic (MCA)
            random rounding (\cref{sec:prism-bench}).
      \item We analyze SR in transformer layers (\cref{sec:analysis}): we establish a
            vector Bienaym\'e--Chebyshev bound that extends prior scalar
            probabilistic forward-error analysis to the full matrix multiplication
            of the MLP projection. We also analyze softmax error propagation by
            deriving an exact decomposition of the cross-entropy loss change
            into signed target shift ($D$) and Kullback--Leibler (KL) distortion, and perform
            a second-order Taylor expansion of the KL term into drift curvature
            ($Q$) and variance penalty ($V$) terms, characterizing the
            site-level trade-offs between SR and RN.
      \item We evaluate on DistilGPT-2 and WikiText-2 (\cref{sec:experiments})
            by lowering $t$ site by site. SR is better in the MLP and RN at the
            language-model head, and the MLP gap widens as more down-projections
            are affected. At $t=6$ in the MLP, SR holds perplexity degradation to
            $15\%$ over the full-precision reference against $121\%$ for RN, and
            comes within $0.3\%$ of the reference by $t=8$.
\end{itemize}

\paragraph{Disclaimer on AI usage.}
AI tools (specifically Anthropic Claude, Google Gemini, and OpenAI Codex) were used during the preparation of this manuscript to assist with prose drafting, editing, and checking the mathematical proofs. All derivations, empirical evaluations, and final text were independently verified by the authors, who take full responsibility for the contents of this work.
